\section{Introduction}
\label{sec:introduction}

Online misinformation is simultaneously a content problem, a social-influence problem, and a measurement problem. False claims can reach large audiences, yet exposure and sharing are not evenly distributed across people or sources \citep{lazer2018,vosoughi2018,grinberg2019,guess2019}. A source may become more visible because it uses emotionally engaging language, familiar visual conventions, or socially reinforced recommendations. However, a source's popularity does not by itself indicate how much false content it contributes. This distinction becomes especially important when a low-credibility source adopts the outward characteristics of an established news organisation. The source may gain attention while retaining its original editorial propensity to publish false material, or it may imitate both presentation and credibility and thereby change the very process that generates false content.

These mechanisms are difficult to isolate with observational platform data. Algorithms, user preferences, social ties, source reach, and content production evolve together. Ethical and practical constraints also prevent researchers from deliberately increasing the attractiveness of a misinformation source in a live platform. Agent-based simulation offers a complementary laboratory in which the mechanisms are explicit, interventions are reproducible, and alternative configurations can be compared while holding other elements constant \citep{macal2010,railsback2019}. Such models do not replace empirical research. Their value lies in making assumptions operational, exposing interactions, and generating conditional claims that can later be calibrated and tested.

This study develops that laboratory using FAKENEWS-ABM, an agent-based model of source evaluation and reposting on social-network sites. The reasoning mechanism is based on endorsement theory \citep{cohen1985}. Users accumulate positive or negative evidence about multiple dimensions of a source, combine that evidence with individual attribute weights, and select a source probabilistically. Word of mouth (WOM) transmits an additional endorsement from a user's contacts. The model extends an earlier endorsement-based line of simulation research on electronic marketplaces, where choices, recognition, and WOM jointly produce market shares \citep{teran2022,teran2024}. In the present application, a source selection represents a repost, and every source independently publishes a true or false item in each period. This makes it possible to distinguish the share of selections received by a target source from the share of selections attached to an item that was actually false.

The research was motivated by a scenario in which a fake-news source copies attributes from a traditional source. Initial experiments treated copying as an undifferentiated operation. Inspection of the model then revealed a consequential dependency: the source-credibility attribute affects both user evaluation and the probability that a source publishes false news. Copying all attributes therefore changes two mechanisms at once. It makes the target resemble the donor in presentation and also makes the target objectively less likely to publish a false item. That scenario remains useful as a mechanistic diagnostic, but it is not a clean representation of strategic camouflage. A more informative treatment copies all attributes except credibility, thereby changing competitive presentation while retaining the target's original false-news propensity.

Memory and WOM add temporal and social dependencies to this distinction. A copied attribute cannot influence a user indefinitely when endorsements expire quickly. An intervention at period 100 may have less opportunity to shape later evidence than an intervention at period 1. WOM can also act in several ways. It may simply reveal an otherwise unknown source, it may penalise a recommender's source after a false publication, or it may reward that source after a true publication. These policies are not interchangeable. A zero outcome is implemented as no outcome endorsement, not as a negative value. Consequently, the model can represent four substantively different rules: penalise or ignore false recommendations, and reward or ignore true recommendations.

We consolidate three experiment phases comprising 2,903 simulation runs. An exploratory phase studies copy scope, memory, timing, network conditions, and common random seeds. A separate baseline phase tests configured false-publication rates and compares no WOM, discovery-only WOM, and truth-sensitive WOM. Finally, strategy, confirmation, and robustness suites examine copied-attribute scope, outcome-specific WOM policies, contact degree, and source reach. The phases are not pooled as if they were identical samples; they are used as complementary levels of evidence. Broad factorial results are treated as exploratory, whereas a smaller set of policy contrasts is repeated with 30 common seeds.

The study addresses four research questions:

\begin{enumerate}
  \item Does the model's baseline reproduce configured false-publication rates, and does truth-sensitive WOM reduce false reposting relative to discovery-only WOM?
  \item How do copied-attribute scope, memory horizon, and activation time affect target-source selection and actual false reposting?
  \item How does the treatment of false and true WOM outcomes alter the effect of non-credibility camouflage?
  \item Are the direction and ordering of the findings robust to contact degree, source reach, and common-seed replication?
\end{enumerate}

The contribution is threefold. First, we provide an operational separation between popularity and realised misinformation in an endorsement-based simulation. Second, we show that a modelling shortcut---using credibility both as a perceived attribute and as an objective content-generation parameter---can reverse the apparent implication of an intervention. Third, we provide a reproducible seeded experiment harness and a validation strategy that distinguishes baseline behaviour, exploratory ranking, and confirmatory contrasts. The resulting claims are about the implemented model, not direct causal estimates for a real platform. Nevertheless, they identify mechanisms and measurement choices that matter for future empirical calibration and for simulation studies of misinformation.

The rest of the paper is organised as follows. Section~\ref{sec:background} reviews misinformation diffusion, social influence, and endorsement-based reasoning. Section~\ref{sec:model} describes the agents, sources, decision process, WOM policy, memory, and scenarios. Section~\ref{sec:design} presents the experimental design and validation procedure. Section~\ref{sec:results} reports baseline, strategy, policy, and robustness results. Section~\ref{sec:discussion} discusses their interpretation and limitations, and Section~\ref{sec:conclusions} concludes.
